\chapter{تمثيلات الزمر المنتهية}\label{ch:b3:representations}

لفهم زمرة مجرّدة، نجعلها تفعل على فضاء متجهي ثم نستعمل الجبر
الخطي --- القيم الذاتية والآثار والجداءات السلَّمية --- على ذلك
الفعل. وهذا البرنامج، أي \emph{نظرية التمثيلات}، فعّال إلى حدّ
مذهل من أجل الزمر المنتهية على $\C$: فكل \omterm{def:b3:representations:rep}{تمثيل} ينشطر إلى تمثيلات
غير قابلة للاختزال (ماشكه)، وتتحدّد التمثيلات غير القابلة
للاختزال بواسطة \emph{طوابعها} (أي آثارها)، وتحقّق الطوابع علاقات
تعامد تجعل الحسابات آلية. ويبني هذا الفصل هذا الحساب وأولى
روائعه --- جداول طوابع الزمر الصغيرة --- وتحصد مسألة نهاية
الأسبوع مبرهنةً تتجاوز بكثير ما تبلغه نظرية الزمر المجرّدة:
مبرهنة بيرنسايد $p^aq^b$. وفي كل ما يلي، $G$ زمرة منتهية وجميع
الفضاءات المتجهية منتهية البُعد على $\C$.

\section{التمثيلات، ماشكه، شور}

\begin{definition}\label{def:b3:representations:rep}
\emph{التمثيل}\index{تمثيل} للزمرة $G$ هو تشاكل $\rho \colon G \to GL(V)$ من أجل فضاء
متجهي $V$ على $\C$؛ ويكون $\dim V$ \emph{درجته}. ويكون فضاء جزئي
$W \subseteq V$ \emph{صامدًا} إذا كان $\rho(g)W \subseteq W$ من أجل كل $g$؛ ويجعل
القصرُ الفضاءَ $W$ \emph{تمثيلًا جزئيًا}. ويكون $\rho$ \emph{غير
قابل للاختزال}\index{تمثيل غير قابل للاختزال} إذا كان $V \neq
0$ ولم
تكن فضاءاته الجزئية الصامدة سوى $0$ و $V$. و\emph{التشاكل} بين
$(\rho, V)$ و $(\sigma, W)$ هو تطبيق خطي $f\colon V \to W$ يحقق $f\rho(g) = \sigma(g)f$ من أجل
كل $g$ (التكافؤ الفعلي)؛ ويُسمّى $f$ التقابلي \emph{تماثلًا}.
\end{definition}

\begin{example}\label{ex:b3:representations:examples}
(a) الدرجة $1$: التشاكلات $G \to \C^\times$.
(b) \emph{التمثيل المنتظم}\index{تمثيل منتظم}: $V = \C^G$ بالأساس
$(e_h)_{h \in G}$، $\rho(g)e_h = e_{gh}$؛ ودرجته $\abs G$.
(c) يعطي فعلٌ تبديلي للزمرة $G$ على مجموعة منتهية $X$ \emph{\omterm{def:b3:representations:rep}{التمثيل}
التبديلي} على $\C^X$: $\rho(g)e_x =
e_{g\cdot x}$.
(d) تفعل $S_n$ على $\C^n$ بتبديل الإحداثيات؛ والفضاء الفائق
$\{\sum x_i = 0\}$ صامد: وهو \emph{\omterm{def:b3:representations:rep}{التمثيل} المعياري}، ودرجته $n - 1$.
\end{example}

\begin{theorem}[ماشكه]\label{thm:b3:representations:maschke}
لكل فضاء جزئي صامد $W$ من \omterm{def:b3:representations:rep}{تمثيل} $(V, \rho)$ متمّمٌ صامد. ومن نتائج
ذلك أن كل \omterm{def:b3:representations:rep}{تمثيل} مجموعٌ مباشر لتمثيلات غير قابلة للاختزال
(\emph{نصف البساطة}).\index{مبرهنة ماشكه}
\end{theorem}

\begin{proof}
ليكن $p \colon V \to V$ \emph{أي} إسقاط صورته $W$ (بأخذ أي متمّم). نأخذ
متوسطه على الزمرة:
\[
\tilde p = \frac1{\abs G}\sum_{g \in G} \rho(g)\,p\,\rho(g)^{-1}.
\]
يرسل كل حد $V$ داخل $W$ (لأن $W$ صامد)، ويُثبّت $W$ نقطةً نقطة:
فمن أجل $w \in W$ لدينا $\rho(g)^{-1}w \in W$، ويُثبّته $p$، ويعيده $\rho(g)$
--- ومنه فإن $\tilde p$ إسقاط على $W$ من جديد. وهو متكافئ فعليًا:
إذ من أجل $h \in G$، يعيد $\rho(h)\tilde p\rho(h)^{-1}$ ترقيم المجموع نفسه. ومنه فإن
$\ker
\tilde p$ متمّم صامد للفضاء $W$. وبتكرار ذلك على الحدود (لأن البُعد
منتهٍ) يتفكّك $V$ إلى تمثيلات غير قابلة للاختزال.
\end{proof}

\begin{theorem}[مبرهنة شور المساعدة]\label{thm:b3:representations:schur}
ليكن $f \colon V \to W$ تشاكلًا بين تمثيلين \emph{غير قابلين للاختزال}.
عندئذٍ يكون $f = 0$ أو يكون $f$ تماثلًا؛ وإذا كان $(V,\rho) = (W,\sigma)$، فإن
$f = \lambda\,\mathrm{id}$ من أجل $\lambda \in \C$ ما. ومنه فإن $\dim\operatorname{Hom}_G(V, W)$ يساوي $1$ إذا
كان $V \cong W$، ويساوي $0$ فيما عدا ذلك.\index{مبرهنة شور
المساعدة}
\end{theorem}

\begin{proof}
الفضاءان $\ker f$ و $\operatorname{im} f$ صامدان (بالتكافؤ الفعلي)، ومنه فكلٌّ
منهما $0$ أو الفضاء كله: أي إما $f = 0$، وإما أن يكون $f$ متباينًا
وصورته الفضاء كله. وإذا كان $V = W$: فللتطبيق $f$ قيمة ذاتية
$\lambda$ (لأن $\C$ مغلق جبريًا)؛ ويكون $f - \lambda\,\mathrm{id}$ تشاكلًا غير متباين
$V \to V$، ومنه فهو $0$. وأما إحصاء البُعد حين $V \cong W$:
فبتثبيت تماثل واحد $u$، يعطي أي تشاكل $f$ التشاكلَ الذاتي $u^{-1}f = \lambda\,\mathrm{id}$:
ومنه $f =
\lambda u$.
\end{proof}

\section{الطوابع والتعامد}

\begin{definition}\label{def:b3:representations:character}
\emph{طابع}\index{طابع} \omterm{def:b3:representations:rep}{التمثيل} $(V, \rho)$ هو $\chi_\rho(g) = \operatorname{tr}\rho(g)$. وهو يحقق
$\chi_\rho(e) = \dim V$، $\chi_\rho(hgh^{-1}) = \chi_\rho(g)$ (لأن الآثار صامدة بالاقتران): أي إن الطوابع
\emph{دوال أصناف}\index{دالة أصناف} --- وهي عناصر الفضاء $\mathcal{CF}(G)$
المؤلَّف من الدوال الثابتة على \omterm{ex:b3:groups:actions}{أصناف الاقتران}، والمزوَّد بالجداء
السلَّمي الهرميتي
\[
\langle \varphi, \psi\rangle = \frac1{\abs G}\sum_{g \in G}
\overline{\varphi(g)}\,\psi(g).
\]
\end{definition}

\begin{proposition}\label{prop:b3:representations:charbasics}
التطبيق $\rho(g)$ قابل للتقطير وقيمه الذاتية جذور للوحدة؛ ولدينا
$\chi_\rho(g^{-1}) = \overline{\chi_\rho(g)}$، و $\abs{\chi_\rho(g)} \leq \chi_\rho(e)$ مع التساوي إذا وفقط إذا كان $\rho(g)$ سلَّميًا.
وتُجمع الطوابع على المجاميع المباشرة: $\chi_{V
\oplus W} = \chi_V + \chi_W$.
\end{proposition}

\begin{proof}
لدينا $\rho(g)^N = \mathrm{id}$ من أجل $N = \abs G$ (لاغرانج): ومنه يُبيد $\rho(g)$
كثيرَ الحدود $X^N - 1$، وهو منشطر بجذور بسيطة، فيكون قابلًا
للتقطير بقيم ذاتية $\lambda_i \in \mu_N$. عندئذٍ $\chi(g^{-1}) = \sum \lambda_i^{-1} = \sum\bar\lambda_i =
\overline{\chi(g)}$، و $\abs{\chi(g)} = \abs{\sum\lambda_i}
\leq \dim V$، مع التساوي
في متراجحة المثلث إذا وفقط إذا كانت جميع القيم $\lambda_i$
متساوية، أي $\rho(g) = \lambda\,\mathrm{id}$. وأما الجمعية على المجاميع المباشرة: فبآثار
الكتل.
\end{proof}

\begin{lemma}\label{lem:b3:representations:homtrace}
ليكن $(V, \rho)$ و $(W, \sigma)$ تمثيلين. المؤثر المتوسِّط على
$\operatorname{Hom}(W, V)$،
\[
c(A) = \frac1{\abs G}\sum_{g}\rho(g)\,A\,\sigma(g)^{-1},
\]
إسقاطٌ على $\operatorname{Hom}_G(W, V)$، ومن أجل التطبيق $\Phi_g \colon A \mapsto \rho(g)A\sigma(g)^{-1}$ لدينا $\operatorname{tr}\Phi_g =
\chi_\rho(g)\,\overline{\chi_\sigma(g)}$.
\end{lemma}

\begin{proof}
التطبيق $c(A)$ متكافئ فعليًا (بإعادة ترقيم المجموع كما في ماشكه)،
و $c$ يُثبّت التطبيقات المتكافئة فعليًا (لأن كل حد يساوي $A$):
ومنه فإن $c$ إسقاط صورته $\operatorname{Hom}_G(W, V)$. وأما الأثر: ففي أُسس، $\Phi_g(A) = BAC$ مع
$B = \rho(g)$ و $C = \sigma(g)^{-1}$؛ وعلى أساس المصفوفات $(E_{kl})$ لدينا
$BE_{kl}C = \sum_{m,n}
b_{mk}c_{ln}E_{mn}$، ومنه فإن معامل $E_{kl}$ في $\Phi_g(E_{kl})$ هو
$b_{kk}c_{ll}$: أي $\operatorname{tr}\Phi_g =
\sum_{k,l}b_{kk}c_{ll} = \operatorname{tr}(B)
\operatorname{tr}(C) = \chi_\rho(g)\chi_\sigma(g^{-1})$، و $\chi_\sigma(g^{-1}) = \overline{\chi_\sigma(g)}$.
\end{proof}

\begin{theorem}[علاقات التعامد الأولى]
\label{thm:b3:representations:orthogonality}
ليكن $\rho, \sigma$ \emph{غير قابلين للاختزال}. عندئذٍ\index{تعامد
الطوابع}
\[
\langle\chi_\sigma, \chi_\rho\rangle =
\begin{cases}
1 & \text{إذا كان } \rho \cong \sigma,\\
0 & \text{فيما عدا ذلك:}
\end{cases}
\]
أي إن الطوابع غير القابلة للاختزال تشكّل عائلة متعامدة ممنظّمة في
$\mathcal{CF}(G)$.
\end{theorem}

\begin{proof}
أثر الإسقاط هو بُعد صورته:
\[
\dim\operatorname{Hom}_G(W, V) = \operatorname{tr} c
= \frac1{\abs G}\sum_g \operatorname{tr}\Phi_g
= \frac1{\abs G}\sum_g \chi_\rho(g)\overline{\chi_\sigma(g)}
= \langle \chi_\sigma, \chi_\rho\rangle,
\]
وتُقيّم مبرهنة شور المساعدة الطرفَ الأيسر بالقيمة $\delta_{\rho \cong
\sigma}$.
\end{proof}

\begin{corollary}\label{cor:b3:representations:multiplicity}
لنفكّك $V \cong \bigoplus_i V_i^{\oplus m_i}$ إلى تمثيلات غير قابلة للاختزال متمايزة ($V_i \not\cong V_j$).
عندئذٍ $m_i = \langle
\chi_{V_i}, \chi_V\rangle$: أي إن الأمثال --- ومنه \omterm{def:b3:representations:rep}{التمثيل} إلى حدود التماثل
--- يحدّدها \omterm{def:b3:representations:character}{الطابع}. وعلاوة على ذلك $\langle \chi_V, \chi_V\rangle = \sum_i
m_i^2$؛ وعلى وجه الخصوص يكون
$V$ \omterm{def:b3:representations:rep}{غير قابل للاختزال} إذا وفقط إذا كان $\langle\chi_V,
\chi_V\rangle = 1$.
\end{corollary}

\begin{proof}
$\chi_V = \sum m_i\chi_{V_i}$
(\cref{prop:b3:representations:charbasics})؛ ثم نأخذ الجداءات السلَّمية مع كل $\chi_{V_i}$
ونستعمل التعامد والتنظيم. وللتمثيلين المتساويَي \omterm{def:b3:representations:character}{الطابع} الأمثالُ
نفسها، ومنه فهما متماثلان.
\end{proof}

\begin{theorem}[التمثيل المنتظم]
\label{thm:b3:representations:regular}
لتكن $\chi_1, \dots, \chi_r$ الطوابعَ المتمايزة غير القابلة للاختزال، ودرجاتها
$n_i = \chi_i(e)$. يتفكّك \omterm{ex:b3:representations:examples}{التمثيل المنتظم} بأمثال $m_i = n_i$؛ ومن ثَمّ
\[
\sum_{i=1}^{r} n_i^2 = \abs G,
\qquad
\sum_i n_i\chi_i(g) = 0 \quad (g \neq e).
\]
\end{theorem}

\begin{proof}
\omterm{def:b3:representations:character}{الطابع} المنتظم: $\chi_{\mathrm{reg}}(g) =
\#\{h : gh = h\}$، وهو يساوي $\abs G$ من أجل $g = e$
ويساوي $0$ فيما عدا ذلك. ومنه $m_i = \langle \chi_i,\chi_{\mathrm{reg}}\rangle =
\frac1{\abs G}\overline{\chi_i(e)}\,\abs G = n_i$. وبتقييم $\chi_{\mathrm{reg}} = \sum n_i\chi_i$ عند $e$
وعند $g \neq e$ نحصل على المتطابقتين المعروضتين.
\end{proof}

\begin{theorem}\label{thm:b3:representations:numberirr}
تشكّل الطوابع غير القابلة للاختزال \emph{أساسًا} متعامدًا ممنظّمًا
للفضاء $\mathcal{CF}(G)$: أي إن عدد التمثيلات غير القابلة للاختزال يساوي
عدد \omterm{ex:b3:groups:actions}{أصناف الاقتران} في $G$.
\end{theorem}

\begin{proof}
لم يبق سوى التمام: لتكن $f \in \mathcal{CF}(G)$ متعامدة مع كل $\chi_i$؛ سنبيّن أن
$f = 0$. من أجل \omterm{def:b3:representations:rep}{تمثيل} $(V,\rho)$، نضع $T_{f,\rho} = \frac{1}{\abs G}
\sum_g \overline{f(g)}\,\rho(g)$. وهو متكافئ فعليًا:
إذ من أجل $h \in
G$،
\[
\rho(h)T_{f,\rho}\rho(h)^{-1} = \frac1{\abs G}\sum_g
\overline{f(g)}\rho(hgh^{-1})
= \frac1{\abs G}\sum_{g'}\overline{f(h^{-1}g'h)}\rho(g') =
T_{f,\rho}
\]
(لأن $f$ \omterm{def:b3:representations:character}{دالة أصناف}). وإذا كان $\rho$ \omterm{def:b3:representations:rep}{غير قابل للاختزال} ودرجته
$n$، أعطت مبرهنة شور أن $T_{f,\rho} = \lambda\,\mathrm{id}$ مع
\[
\lambda = \frac{\operatorname{tr}T_{f,\rho}}{n}
= \frac{1}{n\abs G}\sum_g \overline{f(g)}\,\chi_\rho(g)
= \frac1n\,\langle f, \chi_\rho\rangle = 0 .
\]
إذن $T_{f,\rho} = 0$ على كل \omterm{def:b3:representations:rep}{تمثيل غير قابل للاختزال}، ومنه (على
المجاميع المباشرة) على \emph{كل} \omterm{def:b3:representations:rep}{تمثيل} --- وعلى وجه الخصوص على
\omterm{ex:b3:representations:examples}{التمثيل المنتظم}. ونطبّقه على متجهة الأساس $e_e$: $0 = T_{f,\mathrm{reg}}e_e =
\frac1{\abs G}\sum_g\overline{f(g)}e_g$، وهذا
يفرض أن يكون كل $\overline{f(g)} = 0$. ومنه فإن العائلة المتعامدة الممنظّمة
$(\chi_i)$ تولّد $\mathcal{CF}(G)$، وبُعده عدد \omterm{ex:b3:groups:actions}{أصناف الاقتران}.
\end{proof}

\begin{corollary}[تعامد الأعمدة]
\label{cor:b3:representations:column}
من أجل $g, h \in G$:
\[
\sum_{i=1}^{r}\overline{\chi_i(g)}\,\chi_i(h) =
\begin{cases}
\abs{Z_G(g)} & \text{إذا كان } g, h \text{ مقترنين},\\
0 & \text{فيما عدا ذلك.}
\end{cases}
\]
\end{corollary}

\begin{proof}
لتمثّل $g_1, \dots, g_r$ الأصنافَ، وليكن $c_j$ عدد عناصر الصنف. للمصفوفة
$U_{ij} =
\sqrt{c_j/\abs G}\;\chi_i(g_j)$ من القياس $r \times r$ أسطرٌ متعامدة ممنظّمة
(\cref{thm:b3:representations:orthogonality} مكتوبةً صنفًا صنفًا: $\sum_j \frac{c_j}{\abs G}\chi_i(g_j)\overline{\chi_{i'}(g_j)} =
\delta_{ii'}$)، أي $UU^* = I$؛ ولمصفوفة
مربعة تحقق $UU^* =
I$ يكون أيضًا $U^*U = I$: أي إن الأعمدة متعامدة
ممنظّمة، وهذا يُفكّ إلى المتطابقة المعروضة ($\abs G/c_j = \abs{Z_G(g_j)}$، وهي علاقة
المدار والمثبِّت من أجل الاقتران).
\end{proof}

\begin{proposition}[الطوابع ذات البُعد الواحد؛ الرفع]
\label{prop:b3:representations:onedim}
(a) تكون $G$ أبيلية إذا وفقط إذا كانت درجات جميع تمثيلاتها غير
القابلة للاختزال $1$؛ وعدد الطوابع ذات الدرجة $1$ لأي زمرة $G$
هو $[G : D(G)]$ (وهي طوابع الزمرة الأبيلية المرافقة).
(b) إذا كان $N \trianglelefteq G$، فإن التمثيلات غير القابلة للاختزال للزمرة
$G/N$ تُرفع (بالتركيب مع $G \to G/N$) إلى تلك التمثيلات غير
القابلة للاختزال للزمرة $G$ التي تحتوي نواتها على $N$ بالضبط.
\end{proposition}

\begin{proof}
(a) إذا كانت $G$ أبيلية، كان كل صنف أحاديًا: أي $r = \abs G$،
ويفرض $\sum n_i^2 = \abs G$ أن تكون جميع الدرجات $n_i = 1$؛ وبالعكس، إذا كانت
جميع $n_i = 1$، فإن \omterm{ex:b3:representations:examples}{التمثيل المنتظم} مجموعُ تمثيلات ذات بُعد واحد،
ومنه فإن $\rho_{\mathrm{reg}}(G)$ قابلة للتقطير آنيًا، فتكون تبديلية، ويكون $\rho_{\mathrm{reg}}$
أمينًا: أي إن $G$ أبيلية. والتمثيلات ذات الدرجة $1$ تشاكلات
$G \to \C^\times$ مستقرّها أبيلي: ومنه فهي تتحلّل عبر $G^{\mathrm{ab}} = G/D(G)$
(\cref{exo:b3:groups:9})، ويكون عدد الطوابع المتمايزة للزمرة الأبيلية
$G^{\mathrm{ab}}$ هو $\abs{G^{\mathrm{ab}}}$ (لأن لها ذلك العدد من الأصناف، وجميع درجاتها
$1$). (b) يحفظ التركيب مع الإسقاط عدمَ قابلية الاختزال (لأن
الفضاءات الجزئية الصامدة تتقابل)، ويتحلّل تمثيلٌ تافه على $N$ عبر
\omterm{thm:b3:groups:quotient}{زمرة القسمة} (\cref{thm:b3:groups:firstiso}).
\end{proof}

\section{جداول الطوابع}

\begin{definition}\label{def:b3:representations:table}
\emph{جدول الطوابع}\index{جدول الطوابع} للزمرة $G$ هو المصفوفة
$\bigl(\chi_i(g_j)\bigr)$ من القياس $r \times r$: أسطرها مفهرسة بالطوابع غير
القابلة للاختزال، وأعمدتها \omterm{ex:b3:groups:actions}{بأصناف الاقتران} (مع إظهار أعداد
عناصرها). والأسطر متعامدة ممنظّمة بالنسبة إلى الجداء الموزون،
والأعمدة متعامدة (\cref{cor:b3:representations:column}): فالجدول محدَّد بإفراط شديد،
وهذا ما يجعله قابلًا للحساب.
\end{definition}

\begin{example}[جدول $S_3$]\label{ex:b3:representations:s3}
الأصناف: $e$ (وعدد عناصره 1)، والمبادلات (3)، والدورات الثلاثية
(2)؛ ومنه $r
= 3$ تمثيلًا \omterm{def:b3:representations:rep}{غير قابل للاختزال}، بدرجات $n_i$ تحقق
$\sum n_i^2 = 6$: أي $1, 1,
2$. والدرجة $1$: \omterm{def:b3:representations:character}{الطابع} التافه
$\mathbf 1$ والإشارة $\varepsilon$. ويأتي السطر الأخير من تعامد
الأعمدة (أو من $\chi_{\mathrm{std}} = \chi_{\mathrm{perm}} - \mathbf 1$):
\[
\begin{array}{c|ccc}
S_3 & e & (1\,2)\ [3] & (1\,2\,3)\ [2]\\
\hline
\mathbf 1 & 1 & 1 & 1\\
\varepsilon & 1 & -1 & 1\\
\chi_{\mathrm{std}} & 2 & 0 & -1
\end{array}
\]
تحقق: $\langle\chi_{\mathrm{std}},\chi_{\mathrm{std}}\rangle =
\frac{1}{6}(4 + 0 + 2) = 1$: أي إنه \omterm{def:b3:representations:rep}{غير قابل للاختزال}.
\end{example}

\begin{example}[جدول $S_4$]\label{ex:b3:representations:s4}
الأصناف: $e$ [1]، والمبادلات [6]، والمبادلات المزدوجة [3]،
والدورات الثلاثية [8]، والدورات الرباعية [6]: أي خمسة تمثيلات
غير قابلة للاختزال، $\sum n_i^2 =
24$ مع طابعين من الدرجة $1$ ($\mathbf 1, \varepsilon$؛
$[S_4 :
D(S_4)] = [S_4 : A_4] = 2$): فالدرجات $1, 1, 2, 3, 3$. ويُرفع \omterm{def:b3:representations:character}{الطابع} ذو الدرجة $2$
من $S_4/V \cong S_3$ (\cref{prop:b3:representations:onedim}(b)، حيث $V$ زمرة كلاين)؛ وأما الدرجة
$3$: \omterm{def:b3:representations:rep}{فالتمثيل} المعياري وفتله \omterm{def:b3:representations:character}{بالطابع} $\varepsilon$:
\[
\begin{array}{c|ccccc}
S_4 & e\,[1] & (1\,2)\,[6] & (1\,2)(3\,4)\,[3] &
(1\,2\,3)\,[8] & (1\,2\,3\,4)\,[6]\\
\hline
\mathbf 1 & 1 & 1 & 1 & 1 & 1\\
\varepsilon & 1 & -1 & 1 & 1 & -1\\
\chi_2 & 2 & 0 & 2 & -1 & 0\\
\chi_{\mathrm{std}} & 3 & 1 & -1 & 0 & -1\\
\varepsilon\chi_{\mathrm{std}} & 3 & -1 & -1 & 0 & 1
\end{array}
\]
($\chi_{\mathrm{std}}(g) = \operatorname{fix}(g) - 1$؛ ويقيّم $\chi_2$ جدولَ $S_3$ على صورة كل صنف بترديد
$V$.) وتنجح جميع فحوص تعامد الأسطر والأعمدة --- وتشغيل اثنين
منها هو تمرين التمهيد في \cref{exo:b3:representations:3}.
\end{example}

\begin{method}\label{met:b3:representations:table}
لبناء جدول طوابع: (1) اسرد \omterm{ex:b3:groups:actions}{أصناف الاقتران} وأعداد عناصرها؛ (2)
أحصِ الطوابع ذات الدرجة $1$ بواسطة $G/D(G)$ واكتبها؛ (3) جد
الدرجات الباقية من $\sum n_i^2 = \abs G$ (بتوافيق عددية صغيرة)؛ (4) احصل على
تمثيلات غير قابلة للاختزال بثمن بخس: ارفع من زمر القسمة، واطرح
$\mathbf 1$ من الطوابع التبديلية (وتحقق من $\langle\chi,\chi\rangle = 1$)، واضرب
الطوابع المعروفة في طوابع من الدرجة $1$؛ (5) أنهِ الأسطر المجهولة
بتعامد الأعمدة --- فكل عمود متعامد مع الأعمدة المكتملة أصلًا،
ويحمل عمود $e$ الدرجاتِ. وتحقق من كل شيء بمسح تعامد كامل.
\end{method}

\begin{omfigure}
\begin{tikzpicture}[scale=0.52]
  \draw[thick] (0,0) rectangle (6,6);
  \fill[omDef!25] (0,5) rectangle (1,6);
  \fill[omProp!25] (1,4) rectangle (2,5);
  \fill[omThm!20] (2,0) rectangle (6,4);
  \draw (0,5) rectangle (1,6);
  \draw (1,4) rectangle (2,5);
  \draw (2,0) rectangle (6,4);
  \node[font=\small] at (0.5,5.5) {$1$};
  \node[font=\small] at (1.5,4.5) {$1$};
  \node[font=\small] at (4,2) {$M_2(\C)$};
\end{tikzpicture}

{\small \omterm{ex:b3:representations:examples}{التمثيل المنتظم} للزمرة $S_3$، مقطَّرًا بالكتل: $\C[S_3] \cong \C \times \C \times M_2(\C)$،
بأبعاد $1 + 1 +
4 = 6 = \abs{S_3}$. وفي العموم $\C[G] \cong \prod_i
M_{n_i}(\C)$: فالمتطابقة $\sum n_i^2 = \abs G$ عبارةٌ
عن كتل مصفوفية.}
\end{omfigure}

\section{تمارين}

\begin{exercise}[$\star$]\label{exo:b3:representations:1}
(a) برهن على أن الطوابع غير القابلة للاختزال للزمرة $\Z/n\Z$ هي
$\chi_k(\bar m) = \eu^{2\iu\pi km/n}$، $k = 0, \dots, n-1$، واكتب جدول طوابع $\Z/4\Z$.
(b) تحقق من علاقتَي التعامد كلتيهما عليه --- واعرف المصفوفة: أين
رأى هذا الكتاب مثيلتها من قبل؟
\end{exercise}

\begin{exercise}[$\star$]\label{exo:b3:representations:2}
لتفعل $G$ على مجموعة منتهية $X$ وليكن $\chi$ طابعَ \omterm{def:b3:representations:rep}{التمثيل}
التبديلي $\C^X$.
(a) برهن على $\chi(g) = \abs{\operatorname{Fix}_X(g)}$ و $\langle \mathbf 1, \chi\rangle = \#\{\text{مدارات}\}$ --- أي إن مبرهنة بيرنسايد
المساعدة في الإحصاء (\cref{exo:b3:groups:5}) حسابٌ على الطوابع.
(b) لنفترض أن الفعل متعدٍّ، ومنه $\chi = \mathbf 1 +
\psi$. برهن على أن $\psi$ \omterm{def:b3:representations:rep}{غير
قابل للاختزال} إذا وفقط إذا كان الفعل \emph{متعدّيًا مرتين} (أي
متعدّيًا على الأزواج المرتّبة من نقط متمايزة). \emph{(احسب $\langle\chi,\chi\rangle$
بوصفه عدد المدارات على $X \times X$.)}
(c) اخلص إلى أن \omterm{def:b3:representations:rep}{التمثيل} المعياري للزمرة $S_n$ (حيث $n \geq
2$) \omterm{def:b3:representations:rep}{غير
قابل للاختزال}.
\end{exercise}

\begin{exercise}[$\star$]\label{exo:b3:representations:3}
أعد بناء جدول $S_3$ من الصفر متّبعًا \cref{met:b3:representations:table}، ثم تحقق من
علاقتَي تعامد للأسطر وعلاقتَي تعامد للأعمدة في جدول $S_4$ الوارد
في \cref{ex:b3:representations:s4}. وفكّك \omterm{def:b3:representations:character}{الطابع} التبديلي للزمرة $S_4$ الفاعلة على
$\{1,2,3,4\}$ والطابعَ $\chi_{\mathrm{std}}^2$ (المربّع نقطةً نقطة) إلى تمثيلات
غير قابلة للاختزال.
\end{exercise}

\begin{exercise}[$\star\star$]\label{exo:b3:representations:4}
احسب جدولَي طوابع $D_4$ و $Q_8$. واخلص إلى أن زمرتين غير مماثلتين
يمكن أن يكون لهما جدولا طوابع متطابقان --- فما المعطيات
الزمرية التي يلتقطها الجدول مع ذلك في هذا الزوج (رتب المراكز،
والزمر الأبيلية المرافقة، وعدد العناصر التقابلية)؟ وأيٌّ منها
\emph{يخفق} في التقاطه؟
\end{exercise}

\begin{exercise}[$\star\star$]\label{exo:b3:representations:5}
جدول طوابع $A_4$: الأصناف $e$ [1]، والمبادلات المزدوجة [3]،
و\emph{صنفان} من الدورات الثلاثية [4]، [4].
(a) فسّر انشطار الدورات الثلاثية (قارن \omterm{ex:b3:groups:actions}{المُرَكِّزات} في $S_4$
و $A_4$، كما في \cref{exo:b3:groups:11}).
(b) جد الطوابع الثلاثة ذات الدرجة $1$ (بواسطة $A_4/V \cong
\Z/3\Z$) والطابعَ
ذا الدرجة $3$ (بقصر $\chi_{\mathrm{std}}$ من $S_4$)، وكوِّن الجدول.
(c) اقرأ من الجدول الزمرَ الجزئية الناظمية للزمرة $A_4$ (النوى
$\{g : \chi(g) = \chi(e)\}$ وتقاطعاتها).
\end{exercise}

\begin{exercise}[$\star\star$]\label{exo:b3:representations:6}
(a) برهن على أن $\ker\chi = \{g : \chi(g) = \chi(e)\}$ هي نواة \omterm{def:b3:representations:rep}{التمثيل} الأساسي (\cref{prop:b3:representations:charbasics}،
حالة التساوي).
(b) برهن على أن كل \omterm{def:b3:groups:normal}{زمرة جزئية ناظمية} من $G$ هي تقاطع نوى طوابع
غير قابلة للاختزال. \emph{(مثّل $G/N$ تمثيلًا أمينًا: بتمثيلها
المنتظم.)}
(c) استنتج: تكون $G$ بسيطة إذا وفقط إذا كان $\ker\chi_i = \{e\}$ من أجل كل
$\chi_i$ غير تافه \omterm{def:b3:representations:rep}{وغير قابل للاختزال} --- فالبساطة تُقرأ من \omterm{def:b3:representations:table}{جدول
الطوابع}.
\end{exercise}

\begin{exercise}[$\star\star$]\label{exo:b3:representations:7}
درجات الحالة $\abs G = 8$: برهن على أن نمط درجات زمرة غير أبيلية
رتبتها $8$ هو $(1,1,1,1,2)$، وأن تمثيلها ذا الدرجة $2$ أمين.
وبصورة أعمّ، برهن على أن نمط زمرة غير أبيلية رتبتها $p^3$ هو
$(1^{\,p^2},
p, \dots, p)$، بعدد $p^2$ من الآحاد و $p - 1$ طابعًا من الدرجة $p$.
\emph{(استعمل $[G : D(G)]$ و $\sum n_i^2 = \abs G$؛ وهنا $D(G) = Z(G)$ ورتبته
$p$.)}
\end{exercise}

\begin{exercise}[$\star\star\star$]\label{exo:b3:representations:8}
من أجل زمرتين منتهيتين $G, H$: برهن على أن \omterm{def:b3:representations:character}{دوال الأصناف}
$\chi(g)\psi(h)$ على $G \times H$، من أجل الطوابع $\chi, \psi$
غير القابلة للاختزال للزمرتين $G, H$، هي بالضبط الطوابع غير
القابلة للاختزال للزمرة $G \times H$. \emph{(فالتعامد والتنظيم
حسابٌ مباشر؛ والتمام بإحصاء الأصناف.)} واستنتج جدول طوابع $\Z/2\Z \times \Z/2\Z$
واستنبط من جديد \cref{prop:b3:representations:onedim}(a) من أجل الزمر الأبيلية المنتهية
بواسطة مبرهنة البنية.
\end{exercise}

\begin{exercise}[$\star\star\star$]\label{exo:b3:representations:9}
جدول طوابع $A_5$ (والأصناف بأعداد عناصر $1, 15, 20, 12,
12$ من \cref{exo:b3:groups:11}):
(a) برهن على أن الدرجات هي $1, 3, 3, 4, 5$ \emph{(وهو الحل
الوحيد للمعادلة $\sum n_i^2 = 60$ مع $n_1 = 1$، وباستعمال \cref{exo:b3:representations:6}(c)
مع البساطة، لا يوجد $n_i = 1$ آخر)}.
(b) أنشئ \omterm{def:b3:representations:character}{الطابع} ذا الدرجة $4$ (بالفعل التبديلي على $5$ نقط)
والطابعَ ذا الدرجة $5$ (فالفعل على زمر سيلو الست من النمط $5$
يعطي الدرجة $6 = 1 + 5$؛ وتحقق من عدم قابلية الاختزال)، وأكمل
سطرَي الدرجة $3$ بتعامد الأعمدة: إذ تظهر مدخلات النسبة الذهبية
$\frac{1\pm\sqrt5}2$ على صنفَي الدورات الخماسية.
(c) تحقق من الجدول المكتمل أن $A_5$ بسيطة (\cref{exo:b3:representations:6}(c)).
\end{exercise}

\begin{exercise}[$\star\star$]\label{exo:b3:representations:10}
ليكن $\rho$ تمثيلًا \omterm{def:b3:representations:rep}{غير قابل للاختزال} من الدرجة $n$ وليكن $z
\in Z(G)$.
برهن على $\rho(z) = \lambda_z\,\mathrm{id}$ حيث $\lambda \colon Z(G) \to \C^\times$ تشاكل (ويُسمّى \emph{\omterm{def:b3:representations:character}{الطابع}
المركزي})، واستنتج $\abs{\chi(z)} = n$ من أجل $z$ مركزي. تطبيق: إذا كان
للزمرة $G$ \omterm{def:b3:representations:rep}{تمثيل} أمين \omterm{def:b3:representations:rep}{غير قابل للاختزال}، فإن $Z(G)$ دائرية.
\end{exercise}

\begin{exercise}[$\star\star$]\label{exo:b3:representations:11}
(الإسقاطات المتشابهة النمط) ليكن $(V, \rho)$ تمثيلًا للزمرة $G$
وليكن $\chi_i$ طابعًا \omterm{def:b3:representations:rep}{غير قابل للاختزال} من الدرجة $n_i$. نعرّف
\[
p_i = \frac{n_i}{\abs G}\sum_{g\in G}
\overline{\chi_i(g)}\,\rho(g) \;\in\; \mathcal L(V).
\]
(a) برهن على أن $p_i$ متكافئ فعليًا بالنسبة إلى $G$، واحسب قصره
على \omterm{def:b3:representations:rep}{تمثيل} جزئي \omterm{def:b3:representations:rep}{غير قابل للاختزال} $W \subseteq
V$ طابعه $\chi_j$: فهو
$\delta_{ij}\,
\mathrm{id}_W$ \emph{(بشور؛ ثم خذ الآثار لتعيين السلَّمي)}.
(b) استنتج أن $p_i$ إسقاط على مجموع $V_i$ لجميع التمثيلات
الجزئية غير القابلة للاختزال ذات \omterm{def:b3:representations:character}{الطابع} $\chi_i$ (ويُسمّى
\emph{المركّبة المتشابهة النمط})، وأن $\sum_ip_i =
\mathrm{id}_V$، وأن التفكيك
$V =
\bigoplus_iV_i$ قانوني --- بخلاف الشطر الأدقّ لكل $V_i$ إلى تمثيلات غير
قابلة للاختزال.
(c) من أجل \omterm{ex:b3:representations:examples}{التمثيل المنتظم} للزمرة $S_3$ \omterm{def:b3:representations:character}{وطابع} الإشارة
$\varepsilon$، اكتب $p_\varepsilon$ صراحةً بوصفه عنصرًا من جبر
الزمرة وتحقق يدويًا من $p_\varepsilon^2
= p_\varepsilon$.
\end{exercise}

\begin{exercise}[$\star\star$]\label{exo:b3:representations:12}
(قراءة جدول) جدول طوابع زمرة ما $G$ رتبتها $24$ معروف جزئيًا:
فلها $5$ أصناف بأعداد عناصر $1, 6, 8, 6, 3$، ودرجات
$1, 1, 2, 3, 3$.
(a) استعد الجدول كاملًا: الطابعان الخطيان (أحدهما تافه؛ والآخر
يأخذ القيمة $-1$ على صنفَي العددين $6$ و $6$ بالضبط)، ثم \omterm{def:b3:representations:character}{الطابع}
ذو الدرجة $2$ بتعامد الأعمدة مع عمود المحايد، ثم الطابعان ذوا
الدرجة $3$ بالطريقة نفسها (وأحدهما $\chi_2\chi_4$).
(b) عيّن $G$ (وهي $\cong S_4$: بمقارنة الأصناف بأنماط الدورات)،
واستخرج من الجدول الزمرَ الجزئية الناظمية بواسطة \cref{exo:b3:representations:6}:
نواتا $\chi_2$ (ودليلها $2$: أي $A_4$) و $\chi_3$ (زمرة كلاين
$V_4$)، ولا شيء غيرهما سوى $\{e\}, G$.
(c) اشرح كيف يُظهر الجدول أن $G/V_4 \cong S_3$ \emph{(أي الطوابع تتحلّل
عبر \omterm{thm:b3:groups:quotient}{زمرة القسمة}؟)}.
\end{exercise}

\section{مسألة: مبرهنة بيرنسايد \texorpdfstring{$p^aq^b$}{p^a q^b}}

\begin{problem}[{مسألة نهاية الأسبوع --- قابلية حل الزمر ذات
الرتبة $p^aq^b$}]\label{pb:b3:representations:1}
برهن بيرنسايد سنة 1904 على أن كل زمرة رتبتها ذات عاملين أوليين
على الأكثر قابلةٌ للحل --- وهي عبارة عن زمر مجرّدة لم تُعرَف لها
براهين طوال نصف قرن إلا عبر نظرية الطوابع. وتبني هذه المسألة
البرهان كاملًا، جامعةً \cref{ch:b3:groups} (القابلية للحل) و \cref{ch:b3:modules}
(المقاسات المولَّدة بعدد منتهٍ على $\Z$) وهذا الفصل. وفي كل ما
يلي، $\chi_1, \dots, \chi_r$ هي الطوابع غير القابلة للاختزال للزمرة $G$، مع
$n_i = \chi_i(e)$.

\medskip
\noindent\textbf{الجزء الأول --- الأعداد الصحيحة الجبرية.}
\emph{العدد الصحيح \omterm{def:b3:galois:algebraic}{الجبري}} هو جذر لكثير حدود \emph{واحدي} من
$\Z[X]$.
\begin{enumerate}
  \item برهن على أن $\alpha$ عدد صحيح \omterm{def:b3:galois:algebraic}{جبري} إذا وفقط إذا كانت
        الحلقة $\Z[\alpha]$ مقاسًا على $\Z$ \omterm{def:b3:modules:free}{مولَّدًا بعدد منتهٍ
        من العناصر}.
  \item استنتج أن الأعداد الصحيحة الجبرية تشكّل حلقة جزئية من
        $\C$. \emph{(فإذا كان $\Z[\alpha], \Z[\beta]$ مولَّدين بعدد منتهٍ، فكذلك
        $\Z[\alpha, \beta]$، والمقاسات الجزئية من مقاسات على $\Z$ مولَّدة
        بعدد منتهٍ مولَّدةٌ بعدد منتهٍ، حسب \cref{thm:b3:modules:submodule} وحجّة
        تقديم --- أو مباشرةً: \omterm{def:b3:modules:module}{فالمقاس} الجزئي من $\Z^n$ حر ورتبته
        $\leq n$.)}
  \item برهن على أن العدد الصحيح \omterm{def:b3:galois:algebraic}{الجبري} \emph{الناطق} عددٌ صحيح.
        \emph{(بمبرهنة الجذر الناطق.)}
  \item برهن على أن كل قيمة \omterm{def:b3:representations:character}{طابع} $\chi(g)$ عددٌ صحيح \omterm{def:b3:galois:algebraic}{جبري}.
\end{enumerate}

\medskip
\noindent\textbf{الجزء الثاني --- علاقات مجاميع الأصناف.}
لنثبّت تمثيلًا $(V, \rho)$ \omterm{def:b3:representations:rep}{غير قابل للاختزال} من الدرجة $n$
وطابعه $\chi$. ومن أجل \omterm{ex:b3:groups:actions}{صنف اقتران} $C$، ليكن $S_C = \sum_{g \in C}\rho(g) \in
\mathcal L(V)$.
\begin{enumerate}[resume]
  \item برهن على أن $S_C$ متكافئ فعليًا، ومنه $S_C =
        \omega_C\,\mathrm{id}$ حيث
        \[
        \omega_C = \frac{\abs C\,\chi(g_C)}{n}
        \qquad (g_C \in C \text{ أي ممثل}).
        \]
  \item برهن على $S_CS_{C'} = \sum_{C''}
        a_{CC'C''}\,S_{C''}$ حيث يحصي $a_{CC'C''} \in \N$، من أجل $z \in C''$
        مثبَّت، الأزواجَ $(x, y) \in C \times
        C'$ التي تحقق $xy = z$. واستنتج أن
        الأعداد $\omega_C$ تحقق $\omega_C\,\omega_{C'} = \sum_{C''}
        a_{CC'C''}\,\omega_{C''}$.
  \item اخلص إلى أن كل $\omega_C$ عدد صحيح \omterm{def:b3:galois:algebraic}{جبري}. \emph{(\omterm{def:b3:modules:module}{فالمقاس}
        على $\Z$ المولَّد بالعنصر $1$ وبالأعداد $\omega_C$ حلقةٌ
        مولَّدة بعدد منتهٍ؛ طبّق محك السؤال 1 --- وبدقة أكبر
        برهن على أن $M = \Z\text{-المولَّد }
        (1, (\omega_C)_C)$ يحقق $\omega_{C_0} M \subseteq
        M$ واستعمل حيلة
        المحدد/كايلي--هاملتون، أو حجّة الحلقة الجزئية في السؤال
        2.)}
  \item استنتج \emph{قابلية القسمة لفروبينيوس}: يقسم $n_i$
        العددَ $\abs G$ من أجل كل درجة غير قابلة للاختزال.
        \emph{(احسب $\frac{\abs G}{n} = \frac{\abs G}{n}
        \langle\chi,\chi\rangle = \sum_C \omega_C\,
        \overline{\chi(g_C)}$: فهو عدد صحيح \omterm{def:b3:galois:algebraic}{جبري} وناطق.)}
\end{enumerate}

\medskip
\noindent\textbf{الجزء الثالث --- محك بيرنسايد للبساطة.}
\begin{enumerate}[resume]
  \item ليكن $\chi$ \omterm{def:b3:representations:rep}{غير قابل للاختزال} من الدرجة $n$ وليكن $C$
        صنفًا يحقق $\gcd(\abs C, n) = 1$. باستعمال بيزو والأسئلة من 4 إلى 7،
        برهن على أن $\frac{\chi(g_C)}{n}$ عدد صحيح \omterm{def:b3:galois:algebraic}{جبري}.
  \item لنفترض علاوة على ذلك $0 < \abs{\chi(g_C)} < n$. برهن على استحالة ذلك:
        فللعدد الصحيح \omterm{def:b3:galois:algebraic}{الجبري} $\alpha =
        \chi(g_C)/n$ \emph{مقترنات} --- وهي الأعداد
        $\alpha_\sigma = \frac1n\sigma(\chi(g_C))$ من أجل $\sigma \in \operatorname{Gal}(\Q(\zeta_{\abs G})/\Q)$، وكلٌّ منها متوسط لعدد $n$ من جذور
        الوحدة --- معيار كلٍّ منها $\leq 1$، ومنه يكون الجداء
        $N = \prod_\sigma\alpha_\sigma$ عددًا صحيحًا جبريًا وناطقًا يحقق
        $0 < \abs N < 1$ --- برّر كل عبارة، مستشهدًا بما ورد في \cref{thm:b3:galois:cyclotomicirred}
        من أجل \omterm{def:b3:galois:galois}{زمرة غالوا} وبأن $\sigma$ يبدّل جذور الوحدة. واخلص
        إلى: إما $\chi(g_C) = 0$ وإما $\rho(g_C)$ سلَّمي
        (\cref{prop:b3:representations:charbasics}).
  \item (محك بيرنسايد) ليكن $C \neq \{e\}$ \omterm{ex:b3:groups:actions}{صنف اقتران} عدد
        عناصره \emph{قوة لعدد أولي} $p^k > 1$، ولنفترض أن $G$
        بسيطة وغير أبيلية. يعطي تعامد الأعمدة بين عمود $C$ وعمود
        $e$ أن $1 + \sum_{i \geq 2} n_i\chi_i(g_C) = 0$. برهن على أن طابعًا غير تافه $\chi_i$ مع
        $p \nmid n_i$ يحقق $\chi_i(g_C) \neq 0$ \emph{(وإلا لكان $\frac1p$
        عددًا صحيحًا جبريًا)}؛ ومنه، حسب السؤال 10، يكون
        $\rho_i(g_C)$ سلَّميًا؛ واستنتج تناقضًا مع البساطة
        \emph{(فمجموعة العناصر $g$ التي يكون فيها $\rho_i(g)$
        سلَّميًا زمرةٌ جزئية ناظمية؛ واستعمل الأمانة من
        \cref{exo:b3:representations:6})}. واخلص إلى: \emph{ليس لأي \omterm{def:b3:groups:simple}{زمرة بسيطة} غير
        أبيلية صنفُ اقتران عدد عناصره قوة لعدد أولي وأكبر من
        $1$.}
\end{enumerate}

\medskip
\noindent\textbf{الجزء الرابع --- المبرهنة.}
\begin{enumerate}[resume]
  \item ليكن $\abs G = p^aq^b$ مع $a + b \geq 1$. إذا كانت $G$ بسيطة،
        فبرهن على أنها أبيلية: اختر $z \neq e$ في \omterm{ex:b3:groups:actions}{مركز زمرة} سيلو
        جزئية من النمط $q$ (\cref{thm:b3:groups:pfixed}) وانظر في عدد عناصر صنف
        اقترانه $[G : Z_G(z)]$، وهو قوة للعدد $p$ (ولماذا؟)؛ ثم
        طبّق السؤال 11.
  \item اخلص بالتراجع على $\abs G$: \emph{كل زمرة رتبتها
        $p^aq^b$ \omterm{def:b3:groups:derived}{قابلة للحل}} (بيرنسايد). ولماذا تنهار الحجّة من
        أجل ثلاثة أعداد أولية --- ولماذا يجب أن تنهار، بالنظر
        إلى $\abs{A_5} = 2^2\cdot3\cdot5$؟
\end{enumerate}

\medskip
\noindent\textbf{الجزء الخامس --- جدول طوابع $A_5$.} تستحق أصغرُ
زمرة لا تطالها مبرهنة بيرنسايد صورتَها الكاملة؛ وكل ما يلي لا
يستعمل سوى هذا الفصل مع \cref{exo:b3:groups:11}.
\begin{enumerate}[resume]
  \item ذكّر من \cref{exo:b3:groups:11} \omterm{ex:b3:groups:actions}{بأصناف الاقتران} الخمسة للزمرة $A_5$:
        $\{e\}$، والمبادلات المزدوجة وعددها $15$، والدورات
        الثلاثية وعددها $20$، وصنفان من $12$ دورة خماسية لكلٍّ،
        ممثَّلان بالعنصرين $c =
        (1\,2\,3\,4\,5)$ و $c^2$. وفسّر لماذا تنشطر
        الدورات الخماسية إلى صنفين في $A_5$ رغم أنها تشكّل صنفًا
        واحدًا في $S_5$.
  \item برهن على أن الدرجات غير القابلة للاختزال للزمرة $A_5$ هي
        بالضبط $1, 3, 3, 4, 5$: استعمل $\sum_in_i^2 = 60$ مع $r = 5$
        أصناف، وكون $A_5$ كاملة الاشتقاق ($D(A_5) =
        A_5$)، ومنه فإن
        \omterm{def:b3:representations:character}{الطابع} التافه هو طابعها الخطي الوحيد؛ ثم استبعد كل
        مجموعة متعدّدة أخرى \emph{(اكتب $59$ مجموعَ أربعة مربّعات
        لأعداد صحيحة $\geq 2$: وتحقق من أن ذلك يحدث بطريقة واحدة
        فقط بدرجات معقولة جميعًا)}.
  \item ليكن $\pi$ \omterm{def:b3:representations:character}{الطابع} التبديلي للزمرة $A_5$ على $\{1, \dots, 5\}$:
        $\pi(g) = \#\operatorname{Fix}(g)$، وقيمه $5, 1, 2, 0, 0$ على الأصناف الخمسة.
        احسب $\langle\pi, \mathbf 1\rangle$ و $\langle\pi,
        \pi\rangle$، واستنتج أن $\chi_4 = \pi - \mathbf 1$ \omterm{def:b3:representations:rep}{غير قابل
        للاختزال} ودرجته $4$، وقيمه $4, 0, 1, -1,
        -1$.
  \item اللعبة نفسها على الأزواج غير المرتّبة $\{i, j\}$ وعددها
        $10$: فأعداد النقط الصامدة هي $10, 2, 1, 0, 0$. احسب
        $\langle\pi_{10}, \pi_{10}\rangle$ و $\langle\pi_{10},
        \mathbf 1\rangle$ و $\langle\pi_{10}, \chi_4\rangle$، واستنتج التفكيك $\pi_{10} = \mathbf 1 + \chi_4
        + \chi_5$،
        واحصل على \omterm{def:b3:representations:character}{الطابع} \omterm{def:b3:representations:rep}{غير القابل للاختزال} $\chi_5$ من الدرجة
        $5$ بالقيم $5, 1, -1, 0, 0$.
  \item للطابعين الباقيين غير القابلين للاختزال $\chi_2, \chi_3$
        الدرجةُ $3$. ويحدّد تعامد الأعمدة (كل عمود غير عمود
        المحايد مع عمود المحايد، وكل عمود مع نفسه) قيمهما خارج
        الدورات الخماسية: برهن على $\chi_2(g) = \chi_3(g) = -1$ على المبادلات
        المزدوجة وعلى $0$ على الدورات الثلاثية.
  \item على صنفَي الدورات الخماسية، نضع $x = \chi_2(c)$ و $y =
        \chi_2(c^2)$؛
        ويسمح التناظر بأخذ $\chi_3(c) = y$ و $\chi_3(c^2) = x$. ومن عمود
        $c$ مزدوجًا مع عمود المحايد ومع عمود $c^2$، استنتج
        $x + y = 1$ و $xy = -1$ (وتحقق من القيمة $x^2 +
        y^2 = 3$ التي
        يعطيها عمود $c$ مع نفسه)، ومنه
        \[
        \{x, y\} = \Bigl\{\frac{1 + \sqrt5}2,\ \frac{1 -
        \sqrt5}2\Bigr\} :
        \]
        أي النسبة الذهبية ومقترنها. وكوِّن جدول طوابع $A_5$
        كاملًا.
  \item أجرِ الفحوص: معيار سطر $\chi_2$ يساوي $1$ (استعمل
        $\varphi^2 + \bar\varphi^2 = 3$)، و $\langle\chi_2,
        \chi_3\rangle = 0$، وقابلية القسمة لفروبينيوس (السؤال
        8) من أجل الدرجات الخمس. وأين \emph{ترى} في الجدول
        فرقًا مع $S_5$، التي جميع قيم طوابعها أعداد صحيحة
        ناطقة؟
  \item استنتج من الجدول وحده أن $A_5$ بسيطة: \omterm{def:b3:groups:normal}{فالزمرة الجزئية
        الناظمية} اتحادُ أصناف اقتران يحتوي على $e$ وعدد عناصره
        يقسم $60$ --- تحقق من أن أي مجموع جزئي فعلي للأعداد
        $1 + 15 + 20 + 12 + 12$ يحتوي على الحد $1$ لا يقسم $60$. وقابل ذلك
        بمحك السؤال 11: تحقق من أنه ليس لأي صنف من أصناف
        $A_5$ عددُ عناصر يساوي قوة لعدد أولي وأكبر من $1$.
  \item (خاتمة عشرينية الوجوه) الزمرة $A_5$ هي زمرة دورانات
        المجسّم العشريني الوجوه، والتمثيلان ذوا الدرجة $3$ هما
        الفعلان الهندسيان على $\R^3$. تحقق من متطابقة الأثر:
        فأثر دوران بزاوية $\theta$ هو $1 +
        2\cos\theta$، و $1 + 2\cos\frac{2\pi}5 =
        \frac{1+\sqrt5}2 = \varphi$. وفسّر
        دون أي حساب لماذا يجب أن يحمل \omterm{def:b3:representations:character}{الطابع} الآخر ذو الدرجة $3$
        القيمةَ المقترنة: \omterm{def:b3:galois:galois}{فزمرة غالوا} للامتداد $\Q(\sqrt5)/\Q$
        تفعل على \omterm{def:b3:representations:table}{جدول الطوابع} كله (مدخلةً مدخلةً)، مبدّلةً
        الطوابع غير القابلة للاختزال فيما بينها.
\end{enumerate}

\medskip
\noindent\textbf{الجزء السادس --- تكملات: حدٌّ مركزي والمربّع
التنسوري.}
\begin{enumerate}[resume]
  \item (أدقّ من قابلية قسمة) ليكن $\chi$ \omterm{def:b3:representations:rep}{غير قابل للاختزال} من
        الدرجة $n$. برهن على $\abs{\chi(z)} = n$ من أجل كل $z \in Z(G)$
        \emph{(بمبرهنة شور المساعدة: فإن $\rho(z)$ سلَّمي ورتبته
        منتهية)}، واستنتج من $\langle\chi, \chi\rangle = 1$ الحدَّ
        \[
        n^2 \leq [G : Z(G)] .
        \]
        وبرهن على أن الدرجات غير القابلة للاختزال للزمر غير
        الأبيلية ذات الرتبة $8$ هي $1, 1, 1, 1, 2$ \emph{(خمسة
        أصناف اقتران؛ اكتب $8$ مجموعَ خمسة مربّعات)} وأنها تبلغ
        التساوي $4 = [G : Z(G)]$؛ وتحقق من الحد على $A_5$، التي
        مركزها تافه.
  \item (المربّع التنسوري \omterm{def:b3:representations:character}{للطابع} $\chi_2$) من أجل $g$ ذي رتبة
        منتهية، يكون $\rho(g)$ قابلًا للتقطير بقيم ذاتية هي جذور
        للوحدة؛ واستنتج صيغتَي \omterm{def:b3:representations:character}{الطابع}
        \[
        \chi_{\operatorname{Sym}^2 V}(g)
        = \frac{\chi(g)^2 + \chi(g^2)}2,
        \qquad
        \chi_{\Lambda^2 V}(g)
        = \frac{\chi(g)^2 - \chi(g^2)}2 .
        \]
        وطبّقهما على $\chi_2$ للزمرة $A_5$ \emph{(ولاحظ أن $g^2$
        يمسح صنف $c^2$ حين يمسح $g$ صنف $c$، وبالعكس)}: برهن على
        $\Lambda^2\chi_2 =
        \chi_2$ و $\operatorname{Sym}^2\chi_2 = \mathbf 1 +
        \chi_5$، ومنه
        \[
        \chi_2\otimes\chi_2 = \mathbf 1 + \chi_2 + \chi_5 .
        \]
        وفسّر $\Lambda^2\chi_2 = \chi_2$ هندسيًا بواسطة الجداء المتجهي على $\R^3$.
  \item (تدقيق أخير للجدول) تحقق عدديًا من: تعامد الأعمدة بين
        عمودَي الدورات الخماسية ($\varphi\bar\varphi = -1$)، والقيمة $5 =
        \abs{Z_{A_5}(c)}$ من
        أجل عمود $c$ مع نفسه، وانعدام \omterm{def:b3:representations:character}{الطابع} المنتظم $\sum_i
        n_i\chi_i(g) = 0$ على
        كلٍّ من الأعمدة الأربعة غير عمود المحايد في الجدول.
\end{enumerate}
\end{problem}
